\documentclass[11pt,a4paper]{article}
% Shared preamble for RuPhO English problem statements (match T1 2026 style).
% Use: \input{latex/rupho_en_preamble.tex} after \documentclass.
\usepackage[margin=1in]{geometry}
\usepackage{amsmath,amssymb}
\usepackage{graphicx}
\usepackage{float}
\usepackage{enumitem}
\usepackage{microtype}
\usepackage{caption}
\usepackage{tabularx}
\usepackage{hyperref}

\captionsetup{font=small,labelfont=bf,skip=6pt}

\setlist[itemize]{leftmargin=1.2cm}
\setlist[enumerate]{leftmargin=1.2cm}

% One outer box per subproblem: [id | statement | points] (no inner rules)
\newcommand{\problembox}[3]{%
  \par\addvspace{0.42em}%
  \noindent
  \setlength{\fboxsep}{7.5pt}%
  \setlength{\fboxrule}{0.95pt}%
  \fbox{%
    \begin{minipage}{\dimexpr\linewidth-2\fboxsep-2\fboxrule\relax}
    \begin{tabularx}{\linewidth}{@{}>{\raggedright\arraybackslash\bfseries}p{2.1em} @{\hspace{0.9em}} >{\raggedright\arraybackslash}X @{\hspace{0.9em}} >{\raggedleft\arraybackslash\bfseries}p{2.4em}@{}}
    #1 & #2 & #3
    \end{tabularx}
    \end{minipage}%
  }%
  \par\addvspace{0.32em}%
}


\title{\textbf{Flight to Mars}}
\author{\normalsize Y14 (2014) --- English translation}
\date{}
\begin{document}
\maketitle

A spacecraft with mass $m=10^3~\text{kg}$ is in Earth's orbit around the Sun, far from Earth. To fly into Mars orbit, the vehicle opens its solar panels so that sunlight begins to exert pressure on them. The flight takes place in an ellipse tangent to the circular orbits of the Earth (radius $R_0=1.5 \cdot 10^8~\text{km}$) and Mars (radius $R_1=2.3 \cdot 10^8~\text{km}$)

\problembox{A1}{Find what force $F_0$ must provide the pressure of sunlight on the batteries at the moment the flight begins for the flight to take place.}{10.00}

The batteries are always oriented perpendicular to the sun's rays; neglect the influence of the gravitational fields of the Earth and Mars on the apparatus; the product of the gravitational constant and the mass of the Sun is equal to $GM_0=1.34 \cdot 10^{20}~\text{m}^3/\text{s}^2$. Hint: Light pressure decreases inversely with the square of the distance from the source.

\vspace{1em}
\noindent\textit{Source:} \url{https://pho.rs/p/3985}

\end{document}